\scanpage{144}
% archive sheet number: 143

\[
\left\{\begin{aligned}
\mathrm{Tr}^{\unsure{q|\circ|-}}(u*v)&=bb'\\
\det(u*v)&=\delta(u,v)=b^2c+b^2c'-4cc',
\end{aligned}\right.
\]
\margin{il s'agit ici de la trace et du déterminant en tant qu'élément d'une Alg. quadratique $H\wedge H'$ sur $S$}
% slanted
\margin{\uline{NB} Si $u,v$ sont des \uline{bases}, alors $u*v$ n'est autre que la \unsure{restriction} à $\widetilde S\subset\check V(L_P)$ de la fonction paramétrique $t_{u,v}$ sur $\check V(L_P)$ \unsure{Aff.|Ainsi,} Si on considère $u*v$ comme \deleted{\unsure{la}} \unsure{fct}, \deleted{\illegible{1}} \deleted{\unsure{bil.}} \deleted{\illegible{1}} sur $H^\vee\otimes H'^\vee$, sa restriction à $L_P\subset H^\vee\otimes H'^\vee$ est la fonction paramétrique $t_{uv}$ sur $L_P$.}

Ainsi, la formule \deleted{donnée} plus haut
\deleted{\unsure{donn.}} pour $T_{\varphi,\varphi'}(u,v)$ peut s'écrire

\[
\boxed{T_{\varphi,\varphi'}(u,v)=\varphi\otimes\varphi'\,(u*v)}\,;
\]

Notons aussi les formules

\[
\underset{\text{discr}}{\delta}(u*v)=\delta(u)\delta(v)
\]

ainsi que

\[
\left\{\begin{aligned}
\underbrace{(u*v)'}_{\substack{\text{le $'$ est relatif}\\\text{à l'alg. quad}\\ H\wedge H'\,\deleted{\unsure{,}}}}&=u'*v=u*v'\\
u'*v'&=u*v
\end{aligned}\right.
\]

\uline{NB} Toutes les formules précédentes ont un
sens et sont valables, sans restriction
de car. sur $S,H,H'$, pour \uline{tt} $u\in H$ et $v\in H'$
(base ou pas !)

On récupère $L_{\deleted{\unsure{\varphi}}P}\overset{\text{dfn}}=L_{H,H'}$ à partir de
$\widetilde S=X\wedge X'$ comme l'espace affine enveloppant
(au dessus de $S$) du rev. quad. $\widetilde S$, et
$\delta:L_P\longrightarrow(L\otimes L')^{-1}$ comme l'hom (discriminant !) canonique associé à \deleted{\unsure{le|ce}} revêtement.

